% Λύση του 54ου προβλήματος της ενότητας «Άλυτα Προβλήματα».
\subsection*{Λύση Προβλήματος 54 --- Φαρμακοκινητικό μοντέλο Bateman}

\begin{alist}
  \item Έχουμε
        \[
          C(t)=A\left(e^{-k_et}-e^{-k_at}\right).
        \]
        Παραγωγίζουμε:
        \[
          C'(t)
          =
          A\left(-k_e e^{-k_et}+k_a e^{-k_at}\right).
        \]
        Άρα
        \[
          {
          C'(t)=A\left(k_a e^{-k_at}-k_e e^{-k_et}\right)
          }.
        \]

  \item Για να βρούμε τη χρονική στιγμή μέγιστης συγκέντρωσης, λύνουμε
        την εξίσωση
        \[
          C'(t)=0.
        \]
        Επειδή \(A>0\), έχουμε
        \[
          k_a e^{-k_at}-k_e e^{-k_et}=0.
        \]
        Άρα
        \[
          k_a e^{-k_at}=k_e e^{-k_et}.
        \]
        Διαιρούμε με \(k_e e^{-k_at}>0\):
        \[
          \frac{k_a}{k_e}
          =
          e^{(k_a-k_e)t}.
        \]
        Παίρνουμε λογάριθμο:
        \[
          \ln\left(\frac{k_a}{k_e}\right)
          =
          (k_a-k_e)t.
        \]
        Επειδή \(k_a>k_e\), προκύπτει
        \[
          {
          t_{\max}
          =
          \frac{\ln(k_a/k_e)}{k_a-k_e}
          }.
        \]

        Πριν από αυτή τη στιγμή η απορρόφηση κυριαρχεί και η συγκέντρωση
        αυξάνεται. Μετά από αυτή τη στιγμή η αποβολή κυριαρχεί και η
        συγκέντρωση μειώνεται.

  \item Επειδή
        \[
          k_a>0
          \quad\text{και}\quad
          k_e>0,
        \]
        έχουμε
        \[
          \lim_{t\to+\infty}e^{-k_at}=0
          \quad\text{και}\quad
          \lim_{t\to+\infty}e^{-k_et}=0.
        \]
        Άρα
        \[
          \lim_{t\to+\infty}C(t)
          =
          A(0-0)
          =
          0.
        \]
        Δηλαδή
        \[
          {\lim_{t\to+\infty}C(t)=0}.
        \]
        Φυσικά, αυτό σημαίνει ότι μετά από αρκετό χρόνο το φάρμακο έχει
        ουσιαστικά αποβληθεί από το αίμα.

  \item Για
        \[
          A=10,\qquad k_a=1{,}2,\qquad k_e=0{,}2
        \]
        έχουμε
        \[
          t_{\max}
          =
          \frac{\ln(1{,}2/0{,}2)}{1{,}2-0{,}2}
          =
          \ln6.
        \]
        Άρα
        \[
          {t_{\max}\approx1{,}79\ \text{ώρες}}.
        \]

        Η μέγιστη συγκέντρωση είναι
        \[
          C(t_{\max})
          =
          10\left(e^{-0{,}2t_{\max}}-e^{-1{,}2t_{\max}}\right).
        \]
        Με \(t_{\max}\approx1{,}79\), παίρνουμε
        \[
          C(t_{\max})
          \approx
          10\left(e^{-0{,}358}-e^{-2{,}150}\right).
        \]
        Επομένως
        \[
          C(t_{\max})
          \approx
          5{,}82.
        \]
        Άρα
        \[
          {C_{\max}\approx5{,}82}.
        \]
\end{alist}
